\documentclass[10pt,a4paper]{amsart}
\usepackage[margin=19mm]{geometry}
\usepackage{amsmath,amssymb,mathtools}
\usepackage[scr=rsfs,cal=euler]{mathalfa}
\usepackage{xfrac}
\usepackage[unicode,hidelinks,bookmarks=false]{hyperref}
\newcommand{\ie}{i.e.\ }
\newcommand{\cf}{cf.\ }
\newcommand{\resp}{respectively\ }
\newcommand{\Subsection}{\subsection}
\providecommand{\nobreakdash}{\nobreak\hbox{-}}
\setlength{\emergencystretch}{2em}
\multlinegap0pt
\usepackage{amssymb}
\usepackage{mathtools}
\usepackage{enumerate}
\usepackage{tikz}
\usepackage{caption}
\usepackage[all]{xy}
\newtheorem{lem}{Lemma}
\newcommand{\C}{\mathbb{C}}
\newcommand{\F}{\mathbb{F}}
\newcommand{\Ha}{\mathbb{H}}
\def\Mwgt{\mathrm{Mwgt}}
\title{Secondary cohomology operations and sectional category}
\author{Mark Grant}
\date{Source: arXiv:2512.19461v2 (CC BY 4.0)}
\begin{document}
\maketitle

\noindent\emph{Source and licence.} Secondary cohomology operations and sectional category. Version arXiv:2512.19461v2 (CC BY 4.0), distributed by arXiv under the Creative Commons Attribution 4.0 International licence, http://creativecommons.org/licenses/by/4.0/, as recorded in the arXiv licence metadata for this exact version and shown on the abstract page https://arxiv.org/abs/2512.19461v2. Full licence notice: \texttt{licenses/arxiv-2512.19461-CC-BY-4.0.txt}.\par\smallskip

\noindent\textit{Editorial scope.} Counted unit: the article's lem lem:Mwgttwistor (source0.tex lines 415--417). The article's complete original proof of it is delivered with it (source0.tex lines 419--440, one \texttt{proof} environment of 22 lines). No mathematics is written, repaired or re-derived: the statement and the proof are the article's own lines, in the article's own notation, and no retained line is substituted or re-flowed.  No further source-local context is needed: every cross-reference of every retained line resolves inside the two delivered spans.  Nothing was omitted from any retained span, so the package carries no omission record.  No retained line contains a citation, so the wrapper carries no bibliography and no bibliography entry.  No retained line contains a figure environment, an image-inclusion command or drawing code, so no image is delivered and the package declares no asset dependency.  Editorial formatting only, in addition: an AMS \texttt{amsart} preamble that restates the article's own declarations and macros used by the retained lines, namely the environments \texttt{lem} on separate counters, together with the standard packages those lines need.  The retained environments print in this document as Lemma 1 in order of appearance, whereas the article prints the same items with the numbering of its own full document.  The complete native source of the article as distributed in the arXiv e-print for this exact version is bundled unchanged as \texttt{sources/191480/source0.tex}.
\begin{lem}\label{lem:Mwgttwistor}
For $q$ the twistor bundle, $\Mwgt(q;\F_2)=1$.
\end{lem}

\begin{proof}
Let $H^*(-):=H^*(-;\F_2)$. We have
  \[
 H^*(\Ha P^2)\cong \F_2[a]/(a^3),\qquad |a|=4,
 \]
 \[
 H^*(\C P^5) \cong \F_2[b]/(b^6),\qquad |b|=2,
 \]
 and $q^*:H^*(\Ha P^2)\to H^*(\C P^5)$ is given by $q^*(a)=b^2=Sq^2(b)$.
 
 If $q^*$ were to admit a retraction $s^*:H^*(\C P^5)\to H^*(\Ha P^2)$ as $\mathcal{A}_2$-modules, we would have
 \[
 a=s^* q^*(a)=s^*(Sq^2(b))=Sq^2(s^*(b))=Sq^2(0)=0,
 \]
 a contradiction. So $\Mwgt(q;\F_2)\ge 1$. 
 
 On the other hand, collapse of the spectral sequence for $q(1)$ at the $E_2$-page gives
 \begin{equation}\label{eq:cohE1}
 H^*(E(1)) \cong H^*(\Ha P^2)\otimes_{\F_2} H^*(S^5) \cong \begin{cases} \F_2 & \mbox{if }*=0,4,5,8,9,13, \\ 0 & \mbox{else.} \end{cases}
 \end{equation}
The image of $q(1)^*$ is identified with $H^*(\Ha P^2)\otimes H^0(S^5)$. There are no non-trivial Steenrod operations from $H^*(\Ha P^2)\otimes H^5(S^5)$ to $H^*(\Ha P^2)\otimes H^0(S^5)$ (note that $Sq^3=Sq^1Sq^2$). We can therefore define an $\mathcal{A}_2$-module retraction $s^*$ of $q(1)^*$ by setting $s^*(q(1)^*(x))=x$ for $x\in H^*(\Ha P^2)$ and declaring $s^*$ to be zero outside the image of $q(1)^*$.  Hence $\Mwgt(q;\F_2)\le1$, and the lemma is proved.
 \end{proof}

\end{document}
